\pdfinfoomitdate=1
\pdftrailerid{}
\pdfsuppressptexinfo=15
\documentclass[11pt,letterpaper]{article}
\usepackage[margin=1in]{geometry}
\usepackage[T1]{fontenc}
\usepackage{lmodern,amsmath,amssymb,amsthm,mathtools,booktabs,microtype}
\usepackage[colorlinks=true,linkcolor=black,citecolor=black,urlcolor=blue]{hyperref}
\usepackage{enumitem}
\setlist{itemsep=3pt,topsep=6pt}
\setlength{\parskip}{5pt}
\setlength{\parindent}{0pt}
\setlength{\emergencystretch}{2em}
\newtheorem{theorem}{Theorem}[section]
\newtheorem{lemma}[theorem]{Lemma}
\newtheorem{proposition}[theorem]{Proposition}
\newtheorem{corollary}[theorem]{Corollary}
\theoremstyle{remark}\newtheorem{remark}[theorem]{Remark}
\newcommand{\C}{\mathbb C}
\newcommand{\Q}{\mathbb Q}
\newcommand{\R}{\mathbb R}
\newcommand{\N}{\mathbb N}
\newcommand{\dd}{\mathrm d}
\newcommand{\eps}{\varepsilon}
\newcommand{\ceil}[1]{\left\lceil#1\right\rceil}
\newcommand{\coeff}[2]{[#1^{#2}]}
\hypersetup{pdftitle={Report 118 Inversion sequence asymptotics from iterated kernels},pdfauthor={}}
\begin{document}
\begin{center}
{\small REPORT 118}\par\vspace{8pt}
{\LARGE\bfseries Inversion sequence asymptotics\par}
\vspace{4pt}
{\LARGE\bfseries from iterated kernels\par}
\vspace{9pt}
{\large Rigorous asymptotics and counting index inverses}\par
\vspace{9pt}
{\small 2 October 2026}\par
{\small Self-contained proof and exact reproducibility companion}
\end{center}
\begin{abstract}
For the inversion-sequence classes counted by OEIS A279544 and A279567,
we prove $a_n\sim C4^n n^{-3/2}$ and
$b_n\sim\widehat C(3+2\sqrt2)^n n^{-3/2}$ with positive,
explicitly characterized constants. These establish the leading laws
stated numerically by Britt and Beaton for their classes 214 and 1509.
Normally convergent M\"obius-orbit sums from their exact catalytic
equations give global continuation, convergent local square-root
expansions, and exact certificates of nonzero amplitudes. The second
class also requires a global nonvanishing denominator argument for
an iterated quotient. Both representations give coefficient asymptotics
to every fixed algebraic order, with explicit recurrences. We derive
Lambert-$W_{-1}$ counting-index inverses and shrinking two-ceiling
brackets for exact integer thresholds. The local expansion is convergent; the large-index
expansion is an all-fixed-order asymptotic statement. No complete
exponentially improved transseries, effective numerical inverse cutoff,
nonalgebraicity theorem, literature-priority claim, or external peer-review
claim is made.
\end{abstract}
\setcounter{tocdepth}{1}
{\small\setlength{\parskip}{0pt}\tableofcontents}
\newpage

\section{Model and main conclusions}
\subsection{The counted objects and source statement}
An inversion sequence of length $n$ is a vector $e=(e_1,\ldots,e_n)$
of integers with $0\le e_i<i$. Let $a_n$ count those vectors for which
there are no indices $i<j<k$ satisfying
\begin{equation}\label{avoid214}
e_j\ge e_k\quad\hbox{and}\quad e_i\ge e_k.
\end{equation}
Set $a_0=1$ and $A(z)=\sum_{n\ge0}a_nz^n$. This is an ordinary generating
function, with no factorial normalization. The sequence is
OEIS A279544~\cite{oeis214}; its initial terms are
\[
1,1,2,4,10,26,73,214,651,2040,6549,21453,71485,241702,827603.
\]
The relation notation is $(-,\ge,\ge)$, where the dash imposes no
condition on $e_i,e_j$. Equivalently, the forbidden reduced patterns
are $000,010,100,110,120,210$.

Britt and Beaton~\cite[Section 3.5]{bb} give an exact auxiliary generating
tree and catalytic equation for this class. Their label ``214'' is the
length-seven count, not a shift in the sequence index. Their Section 4.2,
equation (4.8), predicts $a_n\sim C4^n n^{-3/2}$ with $C\approx0.0549097$;
the section explicitly identifies its asymptotics as numerical and
non-rigorous. We use version 3, dated 29 September 2026. The proof below
starts from the exact enumerative model and establishes this specific
leading law together with stronger fixed-order and inverse statements. Section~\ref{sec1509}
gives the same complete treatment of class 1509, whose growth constant
is $3+2\sqrt2$.

The source places the class among generating functions believed to be
nonalgebraic. Neither that belief nor non-$D$-finiteness is proved here.
Having a local square-root singularity is compatible with either an
algebraic or a nonalgebraic generating function.

\subsection{The result and its scope}
\begin{theorem}\label{main214}
The generating function $A$ has radius of convergence $\rho=1/4$.
It extends holomorphically to
\begin{equation}\label{slit214}
\{z:|z|<R\}\setminus[1/4,R)\qquad(1/4<R<1/2),
\end{equation}
and $1/4$ is its unique singularity on the circle of convergence.
For every fixed integer $K\ge0$,
\begin{equation}\label{asym214}
a_n=C4^n n^{-3/2}\left(1+\sum_{j=1}^{K}\frac{c_j}{n^j}
+O(n^{-K-1})\right),\qquad n\longrightarrow\infty,
\end{equation}
where $C>0$. The constants are uniquely specified by the convergent
orbit representation and explicit recurrences below. Moreover the
minimum counting index $N(Y)=\min\{n\ge0:a_n\ge Y\}$ admits the
Lambert-$W_{-1}$ expansion and two-ceiling brackets of
Theorem~\ref{inverse}.
\end{theorem}
The constants hidden in \eqref{asym214} may depend on $K$.
No assertion of convergence as $K\to\infty$ is part of this theorem.
The positive amplitude is certified without fitting coefficients.

The proof has four separate requirements: identify the formal generating
function; prove normal convergence and global continuation; prove that
the square-root term is genuinely nonzero; and transfer with the correct
remainder before inverting. A kernel discriminant alone meets none of
the last three requirements.

\section{Exact generating tree and kernel cancellation}
\subsection{Auxiliary labels and coefficient extraction}
The auxiliary class of~\cite[Section 3.5]{bb} permits the patterns
$000,010,100$ only when their zero symbols have actual value zero.
Its labels $(p,s)$ are lengths of the initial and possible second run of
zeros. The root has label $(0,0)$. Each occurrence of a label creates the
following children, retaining all multiplicities:
\begin{equation}\label{tree214}
\begin{aligned}
(p,s)&\longmapsto(p+1,s),\\
(p,s)&\longmapsto(p+1,s-1)&& (s>0),\\
(p,0)&\longmapsto(p-\ell,\ell)&& (0\le\ell\le p-1),\\
(p,0)&\longmapsto(p-\ell-1,\ell)&& (0\le\ell\le p-2).
\end{aligned}
\end{equation}
An empty index range contributes no child. The desired objects have
$s=0$ and $p\le2$. This completely specifies an integer state recurrence
for $a_n$, and is source equation (3.31).

Let $B(z,x,y)$ record length and the two label coordinates. Put
$F(z,x)=B(z,x,0)$. Every nonroot label has $p\ge1$, so $F(z,0)=1$.
Write $F_p(z)=[x^p]F(z,x)$. The transitions into $(1,0)$ give
\begin{equation}\label{normalize214}
F_1=z(1+F_1+F_2),\qquad
A=1+F_1+F_2=F_1/z.
\end{equation}
The last quotient is removable at zero.

Summing \eqref{tree214}, or clearing the denominators in source
(3.33), gives the exact catalytic relation
\begin{equation}\label{catalytic214}
0=(y+z)(y-x)-(y-zx-zxy)(y-x)B(z,x,y)
+xz(1+y)F(z,y)-z(y+2xy-x^2)F(z,x).
\end{equation}
Each coefficient of $z$ is a polynomial in $x,y$; thus substitutions
below are justified formally after denominator clearing.

\subsection{The root orbit}
Define
\begin{equation}\label{kernel214}
f_z(x)=\frac{zx}{1-zx},\qquad
P(z,x)=zx^2+(2z-1)x+z.
\end{equation}
Substitute $y=f_z(x)$ in \eqref{catalytic214}. Its $B$ term vanishes,
and elementary algebra gives
\begin{equation}\label{F214}
F(z,f_z(x))=P(z,x)F(z,x)+Q(z,x),\qquad
Q(z,x)=\frac{(1+(1-z)x)(1-z-zx)}{1-zx}.
\end{equation}
Set $F(z,x)=1+xG(z,x)$. Then \eqref{F214} becomes
\begin{equation}\label{G214}
G(z,f_z(x))=a(z,x)G(z,x)+b(z,x),
\end{equation}
where
\begin{equation}\label{ab214}
a(z,x)=\frac{(1-zx)P(z,x)}z,
\qquad b(z,x)=x(1-z-zx),\qquad A(z)=G(z,0)/z.
\end{equation}
The apparent $z^{-1}$ is addressed explicitly when its iterated arguments
are substituted; it causes no unproved analytic assumption.

The small root of $P$ is
\begin{equation}\label{r214}
r(z)=\frac{1-2z-\sqrt{1-4z}}{2z}
=\frac{1-\sqrt{1-4z}}{1+\sqrt{1-4z}}\in z\Q[[z]].
\end{equation}
For an arbitrary seed $x$, the iterates have the exact form
\begin{equation}\label{orbit214}
x_j=f_z^{\,j}(x)=
\frac{z^jx}{1-zx(1-z^j)/(1-z)}\qquad(j\ge0).
\end{equation}
It follows either by induction or by summing the geometric series
in the denominator. Define the finite sums and their intended limit by
\begin{align}
\Phi_M(z,x)&=\sum_{k=0}^{M}b(z,x_k)\prod_{j=k+1}^{M}a(z,x_j),\label{partial214}\\
\Phi(z,x)&=\sum_{k\ge0}b(z,x_k)\prod_{j>k}a(z,x_j).\label{phi214}
\end{align}
An empty product is one. No reciprocal of an $a$ factor is used.
The finite sums satisfy
\begin{equation}\label{partialrec214}
\Phi_0=b(z,x),\qquad
\Phi_M=a(z,x_M)\Phi_{M-1}+b(z,x_M).
\end{equation}

\begin{proposition}\label{formal214}
As an identity of formal power series,
\begin{equation}\label{representation214}
A(z)=\Phi(z,r(z))/z.
\end{equation}
\end{proposition}
\begin{proof}
Use seed $x=r(z)$. The initial factor $a(z,r)=0$, so repeated use of
\eqref{G214} gives $G(z,x_{M+1})=\Phi_M(z,r)$. Formula
\eqref{orbit214} gives $x_j=O(z^{j+1})$. For $j\ge1$,
$a(z,x_j)=1+O(z^j)$, while $b(z,x_j)=O(z^{j+1})$.
All products and sums therefore stabilize coefficientwise, and
$G(z,x_{M+1})\to G(z,0)$. Now use \eqref{ab214}.
In particular, this argument never assumes that the original catalytic
series $F(z,x)$ converges at the eventual critical seed $x=1$.
\end{proof}

\section{Normal convergence and dominant singularity control}
\subsection{A holomorphic two-variable orbit sum}
\begin{lemma}\label{normal214}
Choose $0<R<1/2$ and $0<X<(1-R)/R$. The series \eqref{phi214}
and its products converge normally on $|z|\le R$, $|x|\le X$.
They define a jointly holomorphic function near each compact set
with strict room in these inequalities.
\end{lemma}
\begin{proof}
Put $\delta=1-RX/(1-R)>0$. The orbit denominator in
\eqref{orbit214} has modulus at least $\delta$, since
$|z+\cdots+z^j|\le R/(1-R)$. Consequently
\[
|x_j|\le (X/\delta)|z|^j.
\]
For $j\ge1$, $x_j/z$ is holomorphic at zero and uniformly bounded
by $(X/\delta)R^{j-1}$. Expand the first factor as
\begin{equation}\label{expand_a214}
a(z,u)=1+(2-z-1/z)u+(2-2z)u^2-zu^3.
\end{equation}
There are constants $K,K'$ independent of $j,k$ such that
\[
|a(z,x_j)-1|\le KR^{j-1}\quad(j\ge1),\qquad
|b(z,x_k)|\le K'R^k\quad(k\ge0).
\]
The infinite products converge normally and are bounded by
$\exp(K/(1-R))$. The outer sum is therefore dominated by a geometric
series. Uniform limits on slightly larger polydiscs give joint
holomorphy and termwise differentiation on the chosen compact set.
Zeros in finite factors cause no problem because there is no division
by such factors.
\end{proof}
At $z=0$, the first composed factor $a(z,x_1)$ tends to $1-x$, rather
than to one; later factors tend to one. Thus $\Phi(0,x)=x(1-x)$.
This is consistent with the lemma. Since $r(0)=0$, the composition
$\Phi(z,r(z))$ vanishes at zero and its quotient in
\eqref{representation214} is holomorphic there.

\subsection{Continuation on a slit disc}
Use the principal square root on $\C\setminus[1/4,\infty)$.
Then $\operatorname{Re}\sqrt{1-4z}>0$, so the Cayley transform
\eqref{r214} satisfies $|r(z)|<1$. For any $1/4<R<1/2$, one can choose
$1<X<(1-R)/R$. Lemma~\ref{normal214} now proves that
\eqref{representation214} extends the exact germ $A$ to
\eqref{slit214}. In particular, the entire circle $|z|=1/4$ except
$z=1/4$ is crossed holomorphically. This global argument rules out
unexamined competing singularities; inspection of the kernel roots
alone would not do so.

Put $t=\sqrt{1-4z}$ and define
\begin{equation}\label{H214}
H(t)=\frac4{1-t^2}\,
\Phi\left(\frac{1-t^2}{4},\frac{1-t}{1+t}\right)
=\sum_{m\ge0}h_mt^m.
\end{equation}
The lemma applies near $(z,x)=(1/4,1)$, so $H$ is holomorphic for
$|t|$ sufficiently small. This Taylor series is genuinely convergent.
It follows that $A(z)=H(\sqrt{1-4z})$ on the local slit neighborhood,
and differentiation at $t=0$ gives
\begin{equation}\label{h1}
h_1=-8\Phi_x(1/4,1).
\end{equation}
It remains to show that this coefficient is nonzero.

\section{Positive amplitude by exact arithmetic}
\subsection{Rational critical recurrence and a proved tail}
At $z=1/4$ and seed $x=1$, the orbit and its seed derivative are
\begin{equation}\label{criticalorbit}
u_j=\frac3{2\cdot4^j+1},\qquad
v_j^{\mathrm{orb}}=\left.\partial_xx_j\right|_{x=1}
=\frac{9\cdot4^j}{(2\cdot4^j+1)^2}.
\end{equation}
For this subsection write $d_j=v_j^{\mathrm{orb}}$ and
\[
\mathrm a(u)=(1-u/4)(1-u)^2,\qquad
\mathrm b(u)=u(3-u)/4.
\]
Then
\[
\mathrm a'(u)=-\tfrac34(1-u)(3-u),\qquad
\mathrm b'(u)=(3-2u)/4.
\]
If $q_M=\Phi_M(1/4,1)$ and $v_M=\partial_x\Phi_M(1/4,1)$,
\eqref{partialrec214} gives the entirely rational recurrence
\begin{equation}\label{criticalrec}
\begin{aligned}
q_0&=\tfrac12,&v_0&=\tfrac14,\\
q_m&=\mathrm a(u_m)q_{m-1}+\mathrm b(u_m),\\
v_m&=\mathrm a(u_m)v_{m-1}
+d_m\bigl(\mathrm a'(u_m)q_{m-1}+\mathrm b'(u_m)\bigr)
&& (m\ge1).
\end{aligned}
\end{equation}
The derivative uses the old value $q_{m-1}$.

For $0\le u\le1$, one has
\[
0\le\mathrm a\le1,\quad 0\le\mathrm b\le3u/4,\quad
|\mathrm a'|\le9/4,\quad|\mathrm b'|\le3/4,\quad
1-\mathrm a\le9u/4.
\]
For $j\ge1$, $u_j\le(3/2)4^{-j}$ and $d_j\le(9/4)4^{-j}$.
In particular $q_M\le7/8\le9/8$, and the convenient uniform bound
$|\mathrm a'q+\mathrm b'|\le105/32$ yields
\[
|v_M|\le\tfrac14+\tfrac{105}{32}\sum_{j\ge1}d_j
\le\tfrac{347}{128}.
\]
Therefore
\[
|v_m-v_{m-1}|\le\tfrac{347}{128}\tfrac94 u_m
+\tfrac{105}{32}d_m\le\tfrac{16929}{1024}4^{-m}.
\]
Normal convergence identifies the derivative limit, and summation proves
\begin{equation}\label{derivtail}
\left|\Phi_x(1/4,1)-v_M\right|
\le E_M:=\frac{16929}{3072\cdot4^M}.
\end{equation}
At $M=8$, direct exact fraction arithmetic gives
\[
v_8>0.02433,\qquad E_8<0.000085,
\]
so $\Phi_x(1/4,1)>0.024245>0$. The finite inequalities are reproducible
from \eqref{criticalrec}, and no floating-point sign decision is needed.
Consequently
\begin{equation}\label{amplitude}
C=-\frac{h_1}{2\sqrt\pi}
=\frac{4\Phi_x(1/4,1)}{\sqrt\pi}>0.
\end{equation}
Together with \eqref{h1}, this proves that $1/4$ is an actual square-root
singularity. The continuation already established proves the radius and
uniqueness assertions in Theorem~\ref{main214}.

\subsection{Outward rational numerical enclosures}
A certified decimal enclosure for \eqref{amplitude} requires bounds
on $\pi$, not just high-precision evaluation. Use Machin's identity
\[
\pi=16\arctan(1/5)-4\arctan(1/239).
\]
For either arctangent, its alternating rational series is enclosed by
the sum of the first $m$ terms and that sum plus its next term.
Taking $m=100$ and $m=32$, respectively, gives exact rational bounds
$p_-<\pi<p_+$. For $S=10^{80}$ put
\[
s_- =\frac{\lfloor\sqrt{\lfloor S^2p_-\rfloor}\rfloor}{S},
\qquad
s_+ =\frac{\lfloor\sqrt{\lfloor S^2p_+\rfloor}\rfloor+1}{S}.
\]
Integer square roots give $s_-^2\le p_-<\pi<p_+<s_+^2$.
Using \eqref{derivtail} at $M=96$, the exact interval
\[
\frac{4(v_{96}-E_{96})}{s_+}<C<
\frac{4(v_{96}+E_{96})}{s_-}
\]
is narrower than $10^{-56}$. Outward decimal rounding gives
\begin{equation}\label{Cdigits}
\begin{gathered}
0.0549097036253448014962069269284638611865763295943683<C,\\
C<0.0549097036253448014962069269284638611865763295943684.
\end{gathered}
\end{equation}
The interval, rather than any rounded midpoint, is the certified claim.
The matching high-precision numerical constant already appears in
A279544, credited to V\'aclav Kot\v{e}\v{s}ovec on 7 October
2021~\cite{oeis214}. The contribution here is proof and certification
of the amplitude, not discovery of previously unreported digits.

\section{Every fixed coefficient order}
\subsection{Transfer of the convergent local expansion}
For nonintegral $\alpha$, the exact binomial coefficient identity is
\begin{equation}\label{gammaratio}
[z^n](1-4z)^\alpha
=4^n\frac{\Gamma(n-\alpha)}{\Gamma(-\alpha)\Gamma(n+1)}.
\end{equation}
The slit disc provides a standard $\Delta$-domain at $1/4$.
Singularity analysis~\cite[Chapter VI]{fs}, applied to the convergent
Taylor expansion \eqref{H214}, gives, for every fixed $m\ge1$,
\begin{equation}\label{finitegamma}
a_n=4^n\sum_{j=0}^{m-1}
\frac{h_{2j+1}}{\Gamma(-j-1/2)}
\frac{\Gamma(n-j-1/2)}{\Gamma(n+1)}
+O(4^nn^{-m-3/2}).
\end{equation}
For clarity about the remainder, expand $H$ through degree $2m+2$.
The even terms are polynomials in $z$ and have zero coefficients for
all sufficiently large $n$. The first omitted odd term contributes
$O(4^nn^{-m-3/2})$; the remaining Taylor remainder has still smaller
transfer order. No differentiation of an uncontrolled asymptotic
remainder is involved.

Define the standard gamma-ratio polynomials by
\begin{equation}\label{gdef}
\frac{\Gamma(n-\alpha)}{\Gamma(n+1)}
\sim n^{-\alpha-1}\sum_{j\ge0}g_j(\alpha)n^{-j},
\qquad g_0(\alpha)=1.
\end{equation}
Then the exact prescription for the coefficients in \eqref{asym214} is
\begin{equation}\label{Crec}
C_m:=C c_m=\sum_{k=0}^{m}
\frac{h_{2k+1}g_{m-k}(k+1/2)}{\Gamma(-k-1/2)},
\qquad c_0=1.
\end{equation}
In particular, $C_0=C$. For a fully explicit recurrence, let $B_j(x)$
be the Bernoulli polynomials, defined by
$we^{xw}/(e^w-1)=\sum_{j\ge0}B_j(x)w^j/j!$. Set
\begin{equation}\label{grec}
\lambda_j(\alpha)=
\frac{(-1)^{j+1}\bigl(B_{j+1}(-\alpha)-B_{j+1}(1)\bigr)}{j(j+1)},
\qquad
g_m(\alpha)=\frac1m\sum_{j=1}^{m}j\lambda_j(\alpha)g_{m-j}(\alpha).
\end{equation}
This follows by exponentiating the logarithmic Stirling expansion of
the ratio in \eqref{gdef}. It determines any fixed order using rational
arithmetic at the half-integer arguments.

The first two relative corrections are
\begin{equation}\label{c12}
\begin{aligned}
c_1&=\frac38-\frac32\frac{h_3}{h_1},\\
c_2&=\frac{25}{128}-\frac{45}{16}\frac{h_3}{h_1}
+\frac{15}{4}\frac{h_5}{h_1}.
\end{aligned}
\end{equation}
Here $g_1(1/2)=3/8$, $g_1(3/2)=15/8$, and $g_2(1/2)=25/128$.
These identities also provide immediate normalization checks.

\subsection{Certified local jets and the first corrections}
The Taylor coefficients $h_j$ can be computed by rational series
arithmetic in \eqref{partialrec214}, after substituting
$z=(1-t^2)/4$ and $x=(1-t)/(1+t)$. A uniform complex-circle bound
certifies the truncation. On $|t|=1/10$ choose
\[
R=101/400,\qquad X=11/9,\qquad
\delta=1-RX/(1-R),\qquad T=X/\delta.
\]
Then $|z|\le R$, $|x|\le X$ and $|x_j|\le TR^j$. Define
\begin{align*}
A_*&=T+(2+R)TR+(2+2R)T^2R^2+T^3R^3,\\
B_*&=T(1+R+RT).
\end{align*}
Equation \eqref{expand_a214} gives
$|a(z,x_j)-1|\le A_*R^{j-1}$ for $j\ge1$, and
$|b(z,x_j)|\le B_*R^j$. Exact rational inequalities give
\[
A_* /(1-R)<6,\qquad B_*<4,\qquad (1-R)^{-1}<3/2.
\]
Using $e<3$, the finite sums obey
$|\Phi_M|\le B_*e^{A_* /(1-R)}/(1-R)<5000$.
Telescoping \eqref{partialrec214} therefore yields
\[
|\Phi-\Phi_M|
\le\frac{5000A_*+B_*R}{1-R}R^M<40000R^M.
\]
Since $|1/z|\le400/99<5$, if $H_M(t)=\Phi_M(z(t),x(t))/z(t)$, then
\begin{equation}\label{Htail}
\sup_{|t|=1/10}|H-H_M|<200000R^M,
\qquad
|h_j-[t^j]H_M|<200000\cdot10^jR^M.
\end{equation}
The second inequality is Cauchy's estimate. This supplies an effective
certificate for every fixed local coefficient by increasing $M$.

Degree-five rational jets with $M=64$, followed by interval arithmetic
in \eqref{c12}, give the outward enclosures
\begin{equation}\label{correctiondigits}
\begin{aligned}
27.604095556555475840191975&<c_1<27.604095556555475840191976,\\
43.988908237439005751766276&<c_2<43.988908237439005751766277.
\end{aligned}
\end{equation}
For orientation, $h_1\approx-0.19464983128564593535$,
$h_3\approx3.53342590409616985565$, and
$h_5\approx0.37689182228965126927$.
The comparatively large $c_1$ warns against interpreting a short
moderate-$n$ fit as a proof of the leading constant.

\subsection{What is and is not an all-order result}
The series \eqref{H214} converges in a neighborhood of $t=0$.
The expansion \eqref{asym214} means that each fixed truncation has its
own proved algebraic remainder. These are different statements.
No uniform estimate in a truncation order growing with $n$ has been
proved. We do not claim optimal truncation, all subdominant exponential
sectors, Stokes constants, or a complete resurgent transseries. Nor do
the finite coefficient checks replace the analytic continuation and
tail arguments above.

\section{A quotient kernel for class 1509}\label{sec1509}
\subsection{The second counted class and exact equations}
Let $b_n$ count inversion sequences with no $i<j<k$ satisfying
\begin{equation}\label{avoid1509}
e_j\ge e_k\quad\hbox{and}\quad e_i>e_k.
\end{equation}
This is class $(-,\ge,>)$, or avoidance of $100,110,120,210$,
counted by OEIS A279567~\cite{oeis1509}. Put
$\mathcal A(z)=\sum_{n\ge0}b_nz^n$. Its initial terms are
\[
1,1,2,6,21,82,343,1509,6893,32419,156058,765578,3815062.
\]
Source~\cite[Section 3.6, equation (3.34)]{bb} gives the root $(0,0)$
and these multiplicity-preserving transitions:
\begin{equation}\label{tree1509}
\begin{aligned}
(p,s)&\longmapsto(p+1,s),\\
(p,s)&\longmapsto(p+1,s-1)&& (s>0),\\
(p,s)&\longmapsto(p-i,0)&& (s\le1,\ 0\le i<p),\\
(p,s)&\longmapsto(p-\ell,\ell-k)&&
(s\le1,\ 1\le\ell<p,\ 0\le k<\ell).
\end{aligned}
\end{equation}
The target is the sum of all states with $s\le1$. The auxiliary class
allows $100$ when both zero symbols have actual value zero. Let
$\mathcal B(z,x,y)$ be its generating function, and set
\[
J(z,x)=\mathcal B(z,x,0)+[y]\mathcal B(z,x,y),\qquad
\mathcal A(z)=J(z,1).
\]
Write $B_0(x)=\mathcal B(z,x,0)$ and $B_1(x)=[y]\mathcal B(z,x,y)$,
suppressing only the $z$ argument. The source's equations (3.36)--(3.37)
are, equivalently,
\begin{align}
B_0(x)&=1+zxJ(z,x)+\frac{zx}{1-x}\bigl(\mathcal A-J(z,x)\bigr),
\label{B01509}\\
\mathcal B(z,x,y)&=B_0(x)-zxB_1(x)
+\frac{zx(y+1)}y\bigl(\mathcal B(z,x,y)-B_0(x)\bigr)\notag\\
&\quad-\frac{zxy}{(1-x)(x-y)}J(z,x)
+\frac{zxy}{(1-x)(1-y)}\mathcal A
+\frac{zxy}{(1-y)(x-y)}J(z,y).
\label{B1509}
\end{align}
Clear denominators before making formal substitutions. Equation
\eqref{B01509} gives
$B_0(x)=1+zx\mathcal A/(1-x)-zx^2J(z,x)/(1-x)$.
Set $y=f_z(x)=zx/(1-zx)$ in \eqref{B1509}, substitute this expression,
and collect terms. The result is
\begin{equation}\label{J1509}
J(z,f_z(x))=c(z,x)J(z,x)+d(z,x)+e(z,x)\mathcal A(z),
\end{equation}
where, with $K(z,x)=1-x+zx+zx^2$,
\begin{equation}\label{cde1509}
\begin{aligned}
c(z,x)&=\frac{(2zx-1)K(z,x)}{z(x-1)},\\
d(z,x)&=\frac{(2zx-1)(zx+z-1)}{z(zx-1)},\qquad
e(z,x)=-\frac{2x(zx+z-1)}{x-1}.
\end{aligned}
\end{equation}
Thus the scalar reduction follows by rational algebra from the exact
source model, not from a guessed coefficient recurrence.

Define
\begin{equation}\label{rho1509}
\sigma=3-2\sqrt2,\quad \mu=\sigma^{-1}=3+2\sqrt2,\qquad
X(z)=\frac{1-z-\sqrt{1-6z+z^2}}{2z},\quad X(0)=1.
\end{equation}
It satisfies $X=1+zX+zX^2$, so its Taylor coefficients are nonnegative.
Put
\begin{equation}\label{y1509}
y(z)=\frac{X(z)-1}{2}=f_z(X(z)),\qquad
X(\sigma)=1+\sqrt2,\quad y(\sigma)=1/\sqrt2.
\end{equation}
At $x=X(z)$, the kernel $K$ vanishes. The substitution is legitimate
because every coefficient in $z$ of $J$ is a polynomial in $x$;
intermediate divisions by $X-1=2z+O(z^2)$ are interpreted in Laurent
series, or removed by clearing denominators. At that root,
$d=y-1/y$ and $e=1/y$, whence
\begin{equation}\label{boundary1509}
\mathcal A=1-y^2+yJ(z,y).
\end{equation}

\subsection{Iterating in the contracting direction}
Rearrange \eqref{J1509} in the opposite direction from class 214:
\begin{equation}\label{back1509}
J(z,x)=p(z,x)J(z,f_z(x))+q(z,x)+r_*(z,x)\mathcal A(z),
\end{equation}
where
\begin{equation}\label{pqr1509}
\begin{aligned}
p(z,x)&=\frac{z(1-x)}{(1-2zx)K(z,x)},\\
q(z,x)&=\frac{(1-x)(1-z-zx)}{(1-zx)K(z,x)},\\
r_*(z,x)&=\frac{2zx(1-z-zx)}{(1-2zx)K(z,x)}.
\end{aligned}
\end{equation}
For the orbit $x_j=f_z^{\,j}(x)$ of \eqref{orbit214}, put
$P_0=1$, $P_j=\prod_{i=0}^{j-1}p(z,x_i)$, and define
\begin{equation}\label{UV1509}
U(z,x)=\sum_{j\ge0}P_jq(z,x_j),\qquad
V(z,x)=\sum_{j\ge0}P_jr_*(z,x_j).
\end{equation}
Let $U_N,V_N$ stop at $j=N-1$. Finite iteration gives exactly
\[
J(z,x)=U_N(z,x)+V_N(z,x)\mathcal A(z)+P_NJ(z,x_N).
\]
For $x=y(z)=O(z)$, every $p(z,x_j)$ has positive $z$ valuation, so
the last term tends coefficientwise to zero. Combining with
\eqref{boundary1509} proves the formal identity
\begin{equation}\label{F1509}
\mathcal A(z)=\mathcal F(z,y(z)),\qquad
\mathcal F(z,x)=\frac{1-x^2+xU(z,x)}{1-xV(z,x)}.
\end{equation}
Its denominator has constant term one. The essential additional analytic
issue is to exclude zeros of this new denominator throughout the whole
disc, not merely at the positive critical point.

\subsection{A global nonzero denominator certificate}
For $|z|\le\sigma$, positivity and \eqref{rho1509}--\eqref{y1509} give
\[
|z|<R:=43/250,\qquad |X(z)|<M:=483/200,\qquad
|y(z)|<Y:=177/250.
\]
The algebraic kernel factors as
\[
K(z,x)=(1-x/X(z))(1-zX(z)x).
\]
Since $X$ is never zero and $y/X=z/(1-zX)$, at the first orbit point
$x_0=y(z)$ the lower bound is
\begin{equation}\label{k01509}
|K(z,x_0)|\ge k_0:=(1-a)(1-RMY)>0,\qquad
 a=R/(1-RM).
\end{equation}
The exact rational upper bounds
\begin{equation}\label{first1509}
\begin{aligned}
\frac{R(1+Y)}{(1-2RY)k_0}&<\frac45,&
\frac{2RY(1+R+RY)}{(1-2RY)k_0}&<\frac{17}{20},\\
\frac{RY}{1-RY}&<\frac7{50}
\end{aligned}
\end{equation}
imply respectively $|p(z,x_0)|<4/5$, $|r_*(z,x_0)|<17/20$,
and $|x_1|<7/50$. Also
$|q(z,x_0)|\le(1+Y)(1+R+RY)/((1-RY)k_0)$ is finite.

For $|x|\le w:=7/50$, the simpler lower bound
$|K(z,x)|\ge k_s:=1-w-Rw-Rw^2>0$ suffices. Direct rational
inequalities give
\begin{equation}\label{tailbounds1509}
|p|<\tfrac14,\qquad |q|<2,\qquad
|r_*|<\tfrac35|x|,\qquad |f_z(x)|<\tfrac9{50}|x|.
\end{equation}
For example, the upper bounds for $p$ and $r_*/x$ are
$R(1+w)/((1-2Rw)k_s)$ and
$2R(1+R+Rw)/((1-2Rw)k_s)$; the corresponding bounds for $q$ and
$f_z(x)/x$ are $(1+w)(1+R+Rw)/((1-Rw)k_s)$ and $R/(1-Rw)$.
This displays every finite inequality needed for \eqref{tailbounds1509}.

The orbit satisfies $|x_j|\le w(9/50)^{j-1}$ for $j\ge1$.
Normal convergence of both series \eqref{UV1509} follows by geometric
majorants. In particular,
\begin{equation}\label{denom1509}
|yV(z,y)|\le Y\left(\frac{17}{20}
+\frac{(4/5)(3/5)(7/50)}{1-(1/4)(9/50)}\right)<\frac23.
\end{equation}
Hence $|1-yV(z,y)|>1/3$ throughout the closed disc. The factorization
\eqref{k01509} is important: using only the crude first-step triangle
bound would lose the strength required for this global conclusion.
The formal identity \eqref{F1509} is therefore a convergent analytic
representation in $|z|<\sigma$ without a presupposed growth estimate
for the coefficients.

There is also a two-variable neighborhood for $U,V$. On
$|z|\le43/250$, $|x|\le709/1000$, one has
\[
|K(z,x)|\ge1-|x|-R|x|-R|x|^2>0,\qquad |p(z,x)|<5,
\qquad |f_z(x)|<7/50.
\]
The later iterates obey \eqref{tailbounds1509}. All these inequalities
have strict margins, so slightly larger polydiscs still work.
Thus $U,V$ are holomorphic on a neighborhood of this closed polydisc;
$\mathcal F$ is holomorphic near the compact critical curve
$\{(z,y(z)):|z|\le\sigma\}$ by \eqref{denom1509}.

Away from $z=\sigma$ on the boundary circle, the discriminant in
\eqref{rho1509} is nonzero, so $X,y$ extend holomorphically. The
nonzero denominators and contraction estimates persist locally.
At $\sigma$, use a square-root coordinate for $y$ and the just-proved
two-variable holomorphy. Compactness then supplies $\eta>0$ and
$0<\phi<\pi/2$ such that $\mathcal A$ continues to the $\Delta$-domain
\begin{equation}\label{delta1509}
\{z:|z|<\sigma+\eta,\ z\ne\sigma,\ |\arg(z-\sigma)|>\phi\}.
\end{equation}
The next step proves that the possible singularity at $\sigma$ is real.

\subsection{An exact positive critical derivative}
Put $y_*=1/\sqrt2$. For fixed $z=\sigma$, the disc
$|x-y_*|\le1/1000$ lies in $|x|<709/1000$. Its first product factor
has modulus less than five; subsequent factors have modulus less than
$1/4$. Since both subsequent $q$ and $r_*$ have modulus less than two,
the omitted tails in \eqref{UV1509} satisfy, for $N\ge1$,
\begin{equation}\label{UVtail1509}
|U-U_N|,\ |V-V_N|\le E_N:=\frac{40}{3}4^{1-N}.
\end{equation}
Indeed, after the first factor the remaining geometric majorant is
$2/(1-1/4)=8/3$. Cauchy's formula gives at $x=y_*$
\begin{equation}\label{UVderiv1509}
|U_x-(U_N)_x|,\ |V_x-(V_N)_x|\le1000E_N.
\end{equation}

The finite quantities are reproducible with rational interval arithmetic.
Start with zero terminal values and use, backwards for $j=N-1,\ldots,0$,
\[
U_j=q(z,x_j)+p(z,x_j)U_{j+1},\qquad
V_j=r_*(z,x_j)+p(z,x_j)V_{j+1}.
\]
For derivatives, set $\dot x_0=1$,
$\dot x_{j+1}=z\dot x_j/(1-zx_j)^2$, and use, for example,
\[
\dot U_j=q_x(z,x_j)\dot x_j+
 p_x(z,x_j)\dot x_j U_{j+1}+p(z,x_j)\dot U_{j+1}.
\]
Use the identical rule with $q,U$ replaced by $r_*,V$.
All functions being differentiated are the displayed rational functions.
Enclose $\sqrt2$ by scaled integer square roots, then form the intervals
for $\sigma$ and $y_*$. At $N=24$, interval operations on the rational
grid $10^{-60}\mathbb Z$, rounding every operation outward, followed by
\eqref{UVtail1509}--\eqref{UVderiv1509} and quotient differentiation in
\eqref{F1509}, give
\begin{equation}\label{Fx1509}
0.1631959191663560<\mathcal F_x(\sigma,y_*)
<0.1631959201118698.
\end{equation}
In particular, the derivative is strictly positive. This is an exact
finite interval certificate plus a proved infinite tail, not a
floating-point sign test.

Set $t=\sqrt{1-z/\sigma}$ and $z=\sigma(1-t^2)$. The local branch is
\begin{equation}\label{local1509}
y(z)=\frac{1-3z-t\sqrt{1-\sigma z}}{4z}
=y_*-\kappa t+O(t^2),\qquad
\kappa=\frac{\sqrt{1-\sigma^2}}{4\sigma}>0.
\end{equation}
Since $\mathcal F$ is holomorphic near $(\sigma,y_*)$, the expansion
\begin{equation}\label{Blocal1509}
\mathcal F\left(\sigma(1-t^2),
\frac{1-3\sigma(1-t^2)-t\sqrt{1-\sigma^2(1-t^2)}}{4\sigma(1-t^2)}
\right)=\sum_{j\ge0}\widehat h_j t^j
\end{equation}
converges for sufficiently small $|t|$. In particular,
\begin{equation}\label{C1509}
\widehat h_1=-\kappa\mathcal F_x(\sigma,y_*)<0,\qquad
\widehat C=\frac{\kappa\mathcal F_x(\sigma,y_*)}{2\sqrt\pi}>0.
\end{equation}
A rational Machin-series enclosure for $\pi$, combined with the same
outward interval arithmetic, proves
\begin{equation}\label{C1509digits}
0.0660857086342079<\widehat C<0.0660857090170910.
\end{equation}
The more precise numerical value $0.066085708825649431\ldots$ was
already reported by V\'aclav Kot\v{e}\v{s}ovec in A279567 on
7 October 2021~\cite{oeis1509}. The short certificate \eqref{C1509digits} suffices for positivity;
the higher-jet certificate below proves a tighter enclosure.

\subsection{The second all-order theorem}
\begin{theorem}\label{main1509}
For A279567, the radius of convergence is $\sigma=3-2\sqrt2$,
with $\sigma$ the unique dominant singularity. For every fixed $K\ge0$,
\begin{equation}\label{asym1509}
b_n=\widehat C\mu^n n^{-3/2}
\left(1+\sum_{j=1}^{K}\frac{\widehat c_j}{n^j}
+O(n^{-K-1})\right),\qquad \mu=3+2\sqrt2.
\end{equation}
The amplitude is \eqref{C1509}, certified by \eqref{C1509digits}.
Using $g_j$ from \eqref{grec}, all corrections are given by
\begin{equation}\label{Chatrec}
\widehat C\widehat c_m=\sum_{k=0}^{m}
\frac{\widehat h_{2k+1}g_{m-k}(k+1/2)}{\Gamma(-k-1/2)}.
\end{equation}
In particular,
\[
\widehat c_1=\frac38-\frac32\frac{\widehat h_3}{\widehat h_1},\qquad
\widehat c_2=\frac{25}{128}-\frac{45}{16}\frac{\widehat h_3}{\widehat h_1}
+\frac{15}{4}\frac{\widehat h_5}{\widehat h_1}.
\]
The counting-index inverse of Section~\ref{secinverse} applies with
$4,C,c_j$ replaced by $\mu,\widehat C,\widehat c_j$.
\end{theorem}
\begin{proof}
The nonzero coefficient \eqref{C1509} makes $\sigma$ a genuine
square-root singularity. The convergence and continuation leading to
\eqref{delta1509} prove radius and uniqueness. Apply exactly the
transfer argument of \eqref{finitegamma}, replacing $4^n,h_j$ by
$\mu^n,\widehat h_j$. More explicitly, expansion through degree
$2K+4$ has remainder $O(t^{2K+5})$; even powers are polynomials in
$z$, the last retained odd term contributes at most
$O(\mu^nn^{-K-5/2})$, and the remaining terms are smaller after
transfer. The gamma-ratio recurrence proves \eqref{Chatrec} and
\eqref{asym1509}. This also proves source equation (4.10), stated there
as a numerical leading law. The convergent local series permits
certification of each fixed coefficient by finite orbit jets and
Cauchy bounds; no numerical evaluation of a correction is required for
the theorem.
\end{proof}

\subsection{Certified correction jets for class 1509}
The preceding $N=24$ certificate is enough to prove positivity.
A separate degree-five jet calculation gives a much tighter amplitude
and certifies the first two relative corrections. Let
$h=1/10000$ and retain the exact local parametrization
\eqref{local1509}. On $|t|\le h$, subtraction at zero gives
\[
y(t)-y_*=
\frac{t^2-t\sqrt{1-\sigma^2+\sigma^2t^2}}{4\sigma(1-t^2)}.
\]
The radicand has modulus less than one, so
\begin{equation}\label{localbox1509}
|y(t)-y_*|\le\frac{h^2+h}{4\sigma(1-h^2)}<\frac1{5000}.
\end{equation}
Using the rational bounds
$\sigma<171573/10^6$, $y_*<707107/10^6$, and $\sigma>171/1000$,
one obtains $|z|<R$, $|y|<Y$, $|X|=|1+2y|<M$ with the same
$R,M,Y$ used above. The identities $K(z,X)=0$ and $y=f_z(X)$
still hold throughout this $t$ disc, including both signs of $t$.
Therefore the kernel factorization, first-step bound $|p_0|<4/5$,
later bound $|p_j|<1/4$, and $|yV|<2/3$ apply on both branches.
This uses direct complex bounds, not positivity improperly extended
to the second branch.

For $N$ terms the uniform tails now satisfy
\[
|U-U_N|,\ |V-V_N|\le \widetilde E_N
:=\frac{32}{15}4^{1-N}.
\]
The first $q$ bound plus $(4/5)(8/3)$ is less than eight, so
$|U|,|U_N|<8$. For the numerator and denominator of
\eqref{F1509}, write $B=1-y^2+yU$ and $D=1-yV$, and define
$B_N,D_N$ with the finite sums. Then $|B|<8$, $|D|>1/3$, and
$|B-B_N|,|D-D_N|\le Y\widetilde E_N$.
At $N=96$, $Y\widetilde E_N<1/12$, hence $|D_N|>1/4$. Thus
\begin{equation}\label{quotienttail1509}
\left|\frac BD-\frac{B_N}{D_N}\right|
\le\frac{|B-B_N|}{|D_N|}
+\frac{|B|\,|D-D_N|}{|D|\,|D_N|}
<100\widetilde E_N.
\end{equation}
By Cauchy's coefficient estimate,
\begin{equation}\label{jettail1509}
\left|\widehat h_j-[t^j]\mathcal F_N(z(t),y(t))\right|
\le100\widetilde E_N h^{-j}.
\end{equation}
There is no factorial in this coefficient bound.

The finite Taylor jets are evaluated with the same outward rational
interval arithmetic as before. Addition and multiplication are
componentwise sum and truncated convolution. If $A(t)B(t)=1$,
the reciprocal recurrence is
\[
B_0=A_0^{-1},\qquad
B_n=-A_0^{-1}\sum_{k=1}^{n}A_kB_{n-k}.
\]
For the radical in \eqref{local1509}, its degree-four jet is
\[
s+\frac{\sigma^2}{2s}t^2-\frac{\sigma^4}{8s^3}t^4,
\qquad s=\sqrt{1-\sigma^2}.
\]
Multiplication by $t$ means this is sufficient through degree five.
After widening the finite jet coefficients by \eqref{jettail1509},
the formulas of \eqref{C1509} and Theorem~\ref{main1509} give the
following conservative strict enclosures:
\begin{equation}\label{tight1509}
\begin{gathered}
0.066085708825649431003670013119332<\widehat C,\\
\widehat C<0.066085708825649431003670013119333,\\[3pt]
8.0552783200743063565266608644113<\widehat c_1,\\
\widehat c_1<8.0552783200743063565266608644114,\\[3pt]
-17.49290139551774359177529920663<\widehat c_2,\\
\widehat c_2<-17.49290139551774359177529920662.
\end{gathered}
\end{equation}
These bounds certify the two corrections, as well as the indicated
amplitude digits. All higher correction coefficients remain specified
exactly by \eqref{Chatrec}; this report numerically certifies only the
first two relative corrections for each class.

\section{The counting index inverse}\label{secinverse}
\subsection{Leading inverse and formal correction recurrence}
The inverse considered here takes a target count to its required index.
It is not the compositional inverse of the generating function $A(z)$.
The following notation applies to either sequence: let $\mu_0$ be its
growth constant, $C_0$ its positive amplitude, and $c_j^{(0)}$ its
relative corrections. Thus $(\mu_0,C_0,c_j^{(0)})$ equals
$(4,C,c_j)$ for A279544 and $(3+2\sqrt2,\widehat C,\widehat c_j)$
for A279567. In this section abbreviate them to $\mu,C,c_j$,
and write $L=\log\mu$ and $\beta=3/2$. For sufficiently large
real $Y>0$, let
\begin{equation}\label{Wleading}
u(Y)=-\frac{\beta}{L}\,
W_{-1}\left(-\frac{L}{\beta}\left(\frac CY\right)^{1/\beta}\right).
\end{equation}
The branch $W_{-1}$ is the real branch with values at most $-1$ on
$[-1/e,0)$~\cite{lambert}. It selects the large solution of
\begin{equation}\label{leadingeq}
C e^{Lu}u^{-\beta}=Y.
\end{equation}
The principal branch would select the small solution and would be
incorrect for this counting inverse. In logarithmic form,
\begin{equation}\label{logleading}
u(Y)=\frac{\log Y+\beta\log\log Y-\log C-\beta\log L}{L}
+O\left(\frac{\log\log Y}{\log Y}\right).
\end{equation}

Let the logarithmic correction coefficients be defined formally by
\begin{equation}\label{ddef}
\log\left(1+\sum_{j\ge1}c_jw^j\right)=\sum_{j\ge1}d_jw^j.
\end{equation}
Thus $d_1=c_1$, $d_2=c_2-c_1^2/2$, and in general
\begin{equation}\label{drec}
d_m=c_m-\frac1m\sum_{j=1}^{m-1}j d_jc_{m-j}.
\end{equation}
There is a unique formal inverse expansion
\begin{equation}\label{yinverse}
y(Y)=u(Y)+\sum_{j\ge1}e_j u(Y)^{-j},
\end{equation}
obtained by setting all powers of $u^{-1}$ to zero in
\begin{equation}\label{erec}
L(y-u)-\beta\log(y/u)+\sum_{j\ge1}d_jy^{-j}.
\end{equation}
At order $u^{-m}$, the new unknown occurs only as $Le_m$; hence this
is an unambiguous recursion. In particular,
\begin{equation}\label{e123}
\begin{aligned}
e_1&=-d_1/L,\\
e_2&=-\beta d_1/L^2-d_2/L,\\
e_3&=-\beta^2d_1/L^3-\beta d_2/L^2-d_1^2/L^2-d_3/L.
\end{aligned}
\end{equation}
This is the smooth inverse of the asymptotic model. It does not impose
a canonical interpolation on the discrete sequence $a_n$.

\subsection{Strict increase and integer threshold brackets}
Appending the value $n$ to any admissible length-$n$ sequence gives
an admissible length-$(n+1)$ sequence: a new forbidden triple would
have its last, maximal entry at $k=n+1$, which is impossible in
\eqref{avoid214}. The sequence
$(0,1,\ldots,n-1,n-1)$ is also admissible and is not in the image.
For $n=1$ it is $(0,0)$. Therefore $a_{n+1}>a_n$ for $n\ge1$.
For class 1509, appending the maximum gives the same injection,
and the all-zero sequence supplies an extra valid object. Thus
$b_{n+1}>b_n$ for $n\ge1$ as well. In the theorem below, $a_n$
denotes either of these two sequences with its matching constants.

\begin{theorem}\label{inverse}
Let $N(Y)=\min\{n\ge0:a_n\ge Y\}$ and, for a fixed integer $K\ge0$,
put
\[
y_K(Y)=u(Y)+\sum_{j=1}^{K}e_j u(Y)^{-j}.
\]
There are constants $E_K>0$ and $Y_K$ such that, for every $Y\ge Y_K$,
\begin{equation}\label{ceilingbracket}
\ceil{y_K(Y)-E_Ku(Y)^{-K-1}}
\le N(Y)\le
\ceil{y_K(Y)+E_Ku(Y)^{-K-1}}.
\end{equation}
For sufficiently large $Y$, these bounds contain at most two adjacent
integer candidates. The convention at equality is $N(a_n)=n$ for
$n\ge2$.
\end{theorem}
\begin{proof}
The coefficient expansion implies, for every fixed $K$,
\[
\log a_n=Ln-\beta\log n+\log C+
\sum_{j=1}^{K}d_jn^{-j}+O(n^{-K-1}).
\]
Let $F_K(x)=Lx-\beta\log x+\log C+\sum_{j=1}^{K}d_jx^{-j}$.
Choose $D$ and a starting index so that the error at every subsequent
integer is bounded in absolute value by $Dx^{-K-1}$. The smooth
functions $F_K(x)\pm Dx^{-K-1}$ are eventually strictly increasing,
because their derivatives tend to $L>0$.

Invert these two envelopes at $\log Y$. Formal substitution from
\eqref{erec} gives $F_K(y_K)-\log Y=O(u^{-K-1})$; the mean value
theorem, using a fixed positive lower derivative bound, puts both
inverse roots at $y_K+O(u^{-K-1})$. If an integer lies below the
inverse root of the upper envelope, its count is below $Y$. Every
integer at or above the inverse root of the lower envelope has count
at least $Y$. Applying the corresponding ceilings gives
\eqref{ceilingbracket}, after enlarging the fixed constant $E_K$.
Since the real interval has length $2E_Ku^{-K-1}\to0$, its ceilings
are eventually equal or adjacent.
\end{proof}
Already at $K=0$, the real bracket width is $O(1/u)$, not $O(1)$.
It is essential to retain both ceilings: a target may lie arbitrarily
close to an exact sequence value, so unconditional rounding of $y_K$
or an unconditional single ceiling is not justified. The proof gives
existence of $E_K,Y_K$, but this report supplies no effective numerical
values for them. The theorem is therefore not an executable certified
threshold routine for a specified finite target.

\section{Reproducibility and limitations}
The companion package contains this standalone \LaTeX{} source, the
rendered PDF, portable exact checks and fixed data, and deterministic
build and archive scripts. The checks distinguish the 29 displayed
OEIS A279544 prefix terms and 26 displayed A279567 prefix terms from
longer internally generated coefficient lists. Direct brute-force
counting through $n=8$ independently verifies both forbidden-relation
models. For each class, an exact orbit-series construction and a state
recurrence independently agree at all 81 indices $n=0,\ldots,80$.
The class-1509 reconstruction uses the rational orbit quotient. These finite checks identify the models and catch algebraic
or indexing mistakes; they do not prove the asymptotic assertions by
sampling.

The amplitude check uses rational recurrence \eqref{criticalrec}, its
proved infinite tail \eqref{derivtail}, rational arctangent intervals,
and exact integer square roots. The class-214 correction check uses rational
jets and the uniform analytic estimate \eqref{Htail}. The class-1509
checks certify the global denominator bound \eqref{denom1509},
the critical derivative with its Cauchy tail, and the higher-jet
bounds \eqref{quotienttail1509}--\eqref{jettail1509}. Every decimal
interval is rounded outward and checked against its rational
endpoints. The mathematical justification for infinite convergence,
continuation and coefficient transfer is in this report.

The README describes the precise commands, tested toolchain, fixture
schema, checks and failure tests. The package inventory is closed:
changed, missing and unexpected files are rejected. Checksums detect
changes but are not authenticated signatures. The reproducibility
claim is same-toolchain replay in fresh directories, rather than
identity across arbitrary TeX, font or compression versions.

The result establishes the source's class-214 and class-1509 leading
asymptotic laws and all fixed algebraic orders, with their counting-index
inverses in the
stated sense. It does not establish nonalgebraicity or non-$D$-finiteness,
solve every neighboring inversion class, or supply a complete
exponentially improved expansion. The source comparison is explicit
and bounded; no global literature-priority or external peer-review
claim is made.

\section{Further research questions}
The proofs isolate several concrete next questions without presuming
their answers.
\begin{enumerate}[label=\arabic*.]
\item \textbf{Effective inverse thresholds.}
Can one turn the coefficient-transfer remainder into explicit constants
$E_K,Y_K$ in \eqref{ceilingbracket}? This would give a finite-target
certified inverse routine, with exact recurrence checks only at the
remaining one or two candidate indices.
\item \textbf{Algebraic and differential status.}
Do these generating functions fail to be algebraic or $D$-finite, as
suggested in the source? The present local square-root descriptions
do not decide either question.
\item \textbf{Secondary singularities and exponential improvement.}
Can the orbit representations be continued farther, and can their
secondary singularities and exponentially small coefficient
contributions be classified? The bound $R<1/2$ in the class-214
continuation proof is a sufficient domain for this argument, not a
claim of a singularity at $1/2$.
\item \textbf{Other catalytic kernels.}
Which additional inversion-sequence classes admit a contracting
orbit representation with an equally effective nonzero-denominator
certificate? The class-1509 quotient illustrates a new obstruction
that must be handled before transferring a kernel's square root.
\end{enumerate}

\begin{thebibliography}{9}
\bibitem{bb}
N. Britt and N. R. Beaton,
\emph{Completing the enumeration of inversion sequences avoiding
triples of relations}, arXiv:2512.21943v3, 29 September 2026.
Sections 3.5, 3.6 and 4.2, equations (3.31)--(3.37), (4.8) and (4.10).
\url{https://arxiv.org/html/2512.21943v3}.
\bibitem{oeis214}
OEIS Foundation, \emph{A279544}, accessed 2 October 2026.
\url{https://oeis.org/A279544}.
\bibitem{oeis1509}
OEIS Foundation, \emph{A279567}, accessed 2 October 2026.
\url{https://oeis.org/A279567}.
\bibitem{fs}
P. Flajolet and R. Sedgewick, \emph{Analytic Combinatorics},
Cambridge University Press, 2009, Chapter VI.
\url{https://ac.cs.princeton.edu/home/}.
\bibitem{lambert}
R. M. Corless, G. H. Gonnet, D. E. G. Hare, D. J. Jeffrey and D. E. Knuth,
On the Lambert $W$ function, \emph{Advances in Computational Mathematics}
\textbf{5} (1996), 329--359.
\url{https://doi.org/10.1007/BF02124750}.
\end{thebibliography}
\end{document}
